\chapter{Аудит спектральной плотности и масштаба памяти}
\label{ch:v10-nonequilibrium-bath-spectral-memory-audit}

Одна граничная частота не выбирает функцию корреляции. Поэтому проверим
десять кандидатов на полный физический пакет спектральной памяти по шести
условиям:
\begin{enumerate}
\item неотрицательная спектральная плотность;
\item нормированный полный профиль;
\item конечная абсолютная память;
\item совместимость с точными KMS-данными на переходах;
\item выбор существующим родителем независимо от цели;
\item независимое разрушение абсолютной масштабной орбиты.
\end{enumerate}

Проверены экспоненциальный и гауссов профили, омический профиль с мягким
обрезанием, жёсткая зона Бриллюэна, затухающая осцилляция, конечный
спектральный гребень K43, двухванное KMS-завершение, максимум энтропии на
полосе, обратное время часов и наблюдаемое время релаксации.

Точные интегральные свидетели равны
\begin{equation}
 1,\qquad \frac{\sqrt\pi}{2},\qquad
 \frac12,\qquad \frac12
\end{equation}
для экспоненты, гауссиана, затухающего косинуса и компактного треугольного
профиля соответственно. Уже это показывает, что математическая
допустимость не выбирает единственную память.

Матрица кандидатов размера $10\times6$ имеет ранг $6$, но число полных
проходов равно нулю. Лучший результат $5/6$ даёт наблюдаемое время
релаксации: оно разрушает масштабную орбиту, но не выведено существующим
родителем и потому является круговой внешней калибровкой. Все стандартные
аналитические профили достигают не более $4/6$: их форма введена руками.

Кандидаты распадаются на восемь форм профиля и два источника масштаба.
Матрица этого разбиения имеет ранг $2$, а полный вектор происхождения равен
$(0,0)$. Размерная карта
$(\tau_{\rm corr},\omega_{\rm UV},\ell_{\rm cell},v_g)$ сохраняет
ранг/ядро $2/2$.

\begin{theorem}[Полнота отрицательного аудита]
Меню из десяти кандидатов покрывает все шесть критериев независимыми
направлениями, однако ни один кандидат не поставляет одновременно
родительски выбранную форму и некруговой абсолютный масштаб памяти.
\end{theorem}

\begin{nogocriterion}[Профиль нельзя выбрать его пригодностью]
Положительность, нормировка, конечная память и KMS-совместимость допускают
несколько неэквивалентных профилей. Объявление экспоненты или гауссиана
физическим ответом только из-за вычислительного удобства является новым
постулатом, а не выводом.
\end{nogocriterion}

\begin{protocol}\begin{enumerate}
\item проверка кандидатов: \textbf{$10/10$};
\item критериальный ранг: \textbf{$6/6$};
\item полные проходы: \textbf{$0/10$};
\item происхождение формы спектра: \textbf{$0/1$};
\item происхождение абсолютной памяти: \textbf{$0/1$};
\item следующий гейт: \textbf{родитель локального причинного ядра распространения}.
\end{enumerate}\end{protocol}

Машинный сертификат сохранён в
\path{s2t_v10_cell_birth_four_volume_nonequilibrium_bath_spectral_density_memory_scale_candidate_audit_gate_results.json}.